\documentclass[12pt]{article}
\usepackage{amssymb}
\usepackage{amsmath}
\usepackage{amsthm}
\usepackage{mdframed}
\title{prob exercises (stolen)}
\author{Shan Gao}


\begin{document}
\maketitle

\begin{enumerate}
	\item Assume that $S$ is a set and $\{E_n \colon n \geq 1\}$ is a sequence of subsets of $S$. Set $\Sigma:= \sigma(\{E_n \colon n \geq 1\})$. Prove that if $A \in \Sigma$ and $A \subseteq \bigcap_{n \geq 1} E_n$, then $A = \emptyset$ or $A = \bigcap_{n \geq 1} E_n$. 
		\begin{proof}[Solution]
			Using the hint, we consider
			$$
			\Pi := \left\{E \in \Sigma \colon \text{ either } E \supset \bigcap_{n \geq 1} E_n \text{ or } E^c \supset \bigcap_{n \geq 1} E_n \right\}
			$$
			And try to show that $\Pi = \Sigma$, which recovers the solution to our problem(draw a picture). To do this, we need to show that $\Pi$ is a $\sigma$-algebra that contains $E_n$. Trivially, it does contain $\{E_n\}_{n \geq 1}$, so we only need to show that it is a sigma algebra. We need to verify that conditions that make it a sigma algebra.  
			\begin{enumerate}
				\item $S \in \Pi$ and $\emptyset = S\setminus S \in \Pi$ 
				\item Take an arbitrary sequence $\{B_n\colon n \geq 1\} \subseteq \Pi$. Say that for some $n \geq 1$, $B_n \supset \bigcap_{k \geq 1} E_k$, then $\bigcup_{n \geq 1}B_n \supset \bigcap_{k \geq 1} E_k$, so in that case we get closure under countable union. The other case is that for all $n$, $B_n^c \supset \bigcap_{k \geq 1} E_k$, then $\left(\bigcup_{n \geq 1} B_n \right) \supset \bigcap_{k \geq 1} E_k$, which recovers closure under coutable union for \textbf{all} cases.  
			\end{enumerate}
And we are done



		\end{proof}
	\item $S$ is a set, and $\mathcal{C}$ is a collection of subsets of $S$. Consider a set $A$, not necessarily an element of $\mathcal{C}$.
		\begin{itemize}
			\item[i)] Show that $\sigma(\mathcal{C})_{\cap A}$ is a $\sigma$-algebra in $A$. 
				\begin{proof}[solution]
					We verify each condition for a $\sigma$-algebra
					\begin{enumerate}
						\item Obviously $A \in \sigma(\mathcal{C})_{\cap A}$, since $S \in \sigma(\mathcal{C})$ and $A \cap S = A$. 
						\item $\emptyset \in \sigma(\mathcal{C})_{\cap A}$
						\item Take an arbitrary sequence of sets $\{E_n \colon n \geq 1\} \subseteq \sigma(\mathcal{C})_{\cap A}$, each $E_n$ for all $n$, can be represented as an intersection of two sets $D_n \cap A$, where $D_n \in \sigma(\mathcal{C})$, that means 
							\begin{align*}
								\bigcup_{n \geq 1} E_n = \bigcup_{n \geq 1} D_n \cap A = \left(\bigcup_{n \geq 1} D_n\right) \cap A \in \sigma\left(\mathcal{C}\right)_{\cap A} 
							\end{align*}
					\end{enumerate}
				
				\end{proof}
			\item[ii)] Show that $\sigma_A (\mathcal{C}_{\cap A}) = \sigma(\mathcal{C})_{\cap A}$
				\begin{proof}[solution]
					$\subseteq$ is obvious, but let's unwrap the definition
					\begin{align}\label{eq:1}
						\sigma_A(\mathcal{C})_{\cap A} = \sigma_A (\{ B \cap A \colon B \in \mathcal{C})
					\end{align}
					and 
					\begin{align*}
						\sigma(\mathcal{C})_{\cap A} = \{B \cap A \colon B \in \sigma(\mathcal{C})\}
					\end{align*}
					We already know that $\sigma(\mathcal{C})_{\cap A}$ is a sigma algebra containing the set that generated the $\sigma$-alg at \ref{eq:1}. \\
					$\supset$, for the reverse direction, it is a bit less trivial. To do this we consider the hint given. Let 
					\begin{align*}
						\Sigma = \{B \in \sigma(\mathcal{C}) \colon B \cap A \in \sigma_A(\mathcal{C}_{\cap A})\}
					\end{align*}
					We want to show that $\Sigma = \sigma(\mathcal{C})$ and is a $\sigma$-algebra. The proof is the same as exercise i).
				\end{proof}
		\end{itemize}
	\item Prove that $\mathcal{B} ( \mathbb{R} ) = \sigma (\{ (-\infty, x] \colon x \in \mathbb{Q}\}) $ 
		\begin{proof}[solution]
			Since a countable union is countable, to show the $\supset$ direction, we only need to show that any interval $(-\infty, x]$, with $x \in \mathbb{Q}$, can be expressed a countable union of elements in $\mathcal{B} (\mathbb{R})$, Indeed, that is the case with 
			\begin{align*}
				(-\infty, x] = \bigcap_{n \geq 1} \left(-\infty, x + \frac{1}{n}\right) 
			\end{align*}
			Taking the countable union of the above with the complement of each open interval recovers what we wanted. For the reverse inclusion, we can recall from analysis 2 that each open set is a countable union of open intervals, so we only need to show that each $(a,b)$ with $a < b$ and $a , b \in \mathbb{R}$ can be written as a countable union of sets in $\{(-\infty, x] \colon x \in \mathbb{Q}\}$. But first consider than for any $u \in \mathbb{R}$ with $u > a$, $(a,u]$ can be written as $(-\infty, u] \setminus (-\infty, a]$, then
			\begin{align*}
				(a,b) = \bigcup_{n \geq 1} \left ( a, b - \frac{\varepsilon}{n} \right]
			\end{align*}
			letting $\varepsilon = \frac{b-a}{2}$. We can do something quite interesting
				\end{proof}
			\item Show that $\Sigma = \sigma(\{\{x\} \colon x \in \mathbb{R}\}) \subsetneq  \mathcal{B}(\mathbb{R})$
				\begin{proof}[solution]
					Consider the set
					
								\begin{align*}
					\mathcal{G} = \{ E \in \mathbb{R} \colon E \text{ is countable or } E^c \text{ is countable } \} \cup \{\emptyset, \mathbb{R}\}
				\end{align*}
				We want to show that $\mathcal{G} = \Sigma$. We need to show that $\mathcal{G}$ is a $\sigma$-algebra first. The only non trivial condition is closure under countable union. Consider a sequence $\{E_n \colon n \geq 1\} \in \mathcal{G}$, if $A_n$ is countably for all $n$, then we get closure for free. If $A_n^c$ is countable for some $n_0$ , then
				\begin{align*}
					\left(\bigcup_{n \geq 1} A_n \right)^c = \bigcap_{n \geq 1} A_n^c \in \mathcal{G} 
				\end{align*}
				That means, by definition of $\mathcal{G}$, $\bigcup_{n \geq 1} A_n \in \mathcal{G}$, thus it is a sigma algebra. It's obvious that $\Sigma \subseteq \mathcal{G}$, and we've shown the other direction thus $\Sigma = \mathcal{G}$. But then there's a contradiction since we can take any open interval like $(0,1)$, which is in $\mathcal{B}(\mathbb{R})$, but not in $\Sigma$. 
				\end{proof}
			\item \label{q:5}
				Consider the following set 
				\begin{align*}
					\Sigma_0 = \left\{\bigcup_{i=1}^k (a_i, b_i] \colon k \in \mathbb{N} \text{ and } a_1 < b_1 < a_2 < b_2 < \ldots < a_k < b_k \right\}
				\end{align*}
				This is an algebra (not a $\sigma$-algebra). Consider the following set function $\mu \colon \Sigma_0 \to [0,\infty]$
				\begin{align*}
					\text{for all $A \in \Sigma_0$ } \quad \mu(A) = \begin{cases}
						1 & \text{ if $A \supseteq (\frac{1}{2}, \frac{1}{2} + \varepsilon] $ for some $\varepsilon > 0$} \\
						0 & \text{otherwise}
					\end{cases}
				\end{align*}
				Show that $\mu$ is finitely additive but not countably additive.
				\begin{proof}[solution]
					Denote $F_n = \bigcup_{i = 1}^{k_n} (a_i^n, b_i^n]$ to be a finite sequence such that $F_i \cap F_j = \emptyset$ if $i \neq j$. Considering the set $\mathcal{F} = \{F_n \colon 1 \leq n \leq N\}$ for some $n \in \mathbb{N}$. We break it down into cases. \\
					Case 1: $\mu(F_n) = 0$ for all $n \leq N$, that means for all $n$ less than $N$, $F_n \cap (\frac{1}{2}, \frac{1}{2} + \varepsilon_n] = \emptyset$, but that means, if we take the smallest epsilon, call it $\varepsilon_0 = \min_{n \leq N} \varepsilon_n$, then $\bigcup_{n}^N F_n \cap \varepsilon_0 = \emptyset$. \\
					Case 2: There exists some $k \in \mathbb{N}$ such that $\mu(F_k) = 1$, taking the union we just get a bigger set. 
					To show that it's not countably additive, we just need to show that it's not countinuous at $\emptyset$, let $A_n = (\frac{1}{2}, \frac{1}{2} + \frac{1}{n})$, clearly $A_n \downarrow \emptyset$, but $\mu(A_n) \to 1$. 
			\end{proof}
		\item Consider the measure space $(\mathbb{R}, \mathcal{B}(\mathbb{R}), \lambda)$, where $\mathbb{B}(\mathbb{R})$ is the Borel $\sigma$-algebra and $\lambda$ is the Lebesgue measure.
			\begin{itemize}
				\item[i)] Show that, for every $B \in \mathcal{B}(\mathbb{R})$ and $a \in \mathbb{R}$. 
					\begin{align*}
						a+B := \{a+x \colon x \in B\} \in \mathcal{B}(\mathbb{R})
					\end{align*}
					and $\lambda(B) = \lambda(a+B)$, i.e., the Lebesgue measure is translation invariant. 
					\begin{proof}[solution]
						Let
						\begin{align*}
							\Sigma = \{a + B \colon B \in \mathcal{B}(\mathbb{R})\} 
						\end{align*}
						We can show the first statement by proving that $\Sigma = \mathcal{B}(\mathcal{R})$. We show that it's a sigma algebra first. Obviously, $\mathbb{R} \in \Sigma$, and $\emptyset \in \Sigma$. Then let $A \in \Sigma$, that means $A = a + B$ for some $B \in \mathcal{B}(\mathbb{R})$, that means $A^c = (a+B)^c = a + B^c \in \Sigma$. Let a sequence $\{A_n \colon n\geq 1\} \in \Sigma$, then it can decompose into $A_n = B_n + a$ where $B_n \in \mathcal{B}(\mathbb{R})$ for all $n$. Then 
						\begin{equation*}
							\bigcup_{n \geq 1} A_n = \bigcup_{n \geq } B_n + a = \left(\bigcup_{n \geq 1} B_n\right) +a \in \Sigma
						\end{equation*}
						So $\Sigma$ is a $\sigma$-algebra which contains all the open intervals of $\mathbb{R}$, that means, $\Sigma = \mathcal{B}(\mathbb{R})$, we've shown the first statement. Moving on to the second statement. We define a new set function $\mu : \mathcal{B}(\mathbb{R}) \to [0,\infty)$, where $\mu(B) = \lambda(a + B)$, where $B \in \mathcal{B}(\mathbb{R})$. Let $(B_n)_{n \geq 1}$ be a sequence in $\mathcal{B}(\mathbb{R})$, it's easy to see if that if all $B_i$ are disjoint then the translation $a+B_i$ form disjoint sets as well. Then 
						\begin{align*}
							\mu\left(\bigcup_{i\geq 1} B_i\right)&= \lambda\left(\bigcup_{i \geq 1}(B_i) + a \right) \\
							&=\lambda\left(\bigcup_{i \geq 1}B_i + a \right)\\
							&= \sum_{i \geq 1} \lambda(B_i + a)\\
							&= \sum_{i \geq 1} \mu(B_i)
						\end{align*}
						Thus $\mu$ is countably additive and forms a measure on $\mathcal{B}(\mathbb{R})$. To finish off the proof, consider the set of open intervals, which form a $\pi$-system on the reals. It's clear that $\mu = \lambda$ on the latter $\pi$-system, thus they agree on the $\sigma$-algebra generated by that $\pi$-system, therefore $\mu = \lambda$ on $\mathcal{B}(\mathbb{R})$ 
					\end{proof}
					\item Given another non trivial measure $\mu$ on $\mathcal{B}(0,1]$ that is translation invariant, show that $\mu = c \lambda$ for some constant $c$. 
					\begin{proof}[Solution.]
						To do this, we inspire ourselves off the definition of the Lebesgue measure. Consider the interval $(0,1]$, it can be written as a finite union of (letting $n$ be any nature number)
						$$ (0,1] = \bigcup_{k=0}^{n-1} \left( \frac{k}{n}, \frac{k+1}{n}\right]$$
						applying the measure 
						\begin{align*}
							c=\mu(0,1] &= \mu\left(\bigcup_{k=0}^{n-1} \left( \frac{k}{n}, \frac{k+1}{n}\right] \right) \\
							&= \mu\left(\bigcup_{k=0}^{n-1} \left( \frac{k}{n} + \left(0, \frac{1}{n}\right]\right) \right)\\
							&= n \cdot \mu\left(0,\frac{1}{n}\right]
						\end{align*}
						after guessing the constant, we set $\mu\left(0,\frac{1}{n}\right] = \frac{c}{n}$. Similarly, we can apply the same reasoning to get, for any $m \leq n$ and $n$, $\mu\left(0,\frac{m}{n}\right) = \frac{m}{n} \cdot c$. Since, applying translational invariance, 
						\begin{equation*}
							\mu \left(0, \frac{m}{n} \right] = \mu \left(\bigcup_{i=0}^{m-1} \left(\frac{i}{n},\frac{i+1}{n}\right]\right) = m\cdot \mu\left(0,\frac{1}{n}\right] = c \cdot \frac{m}{n}
						\end{equation*}
						By construction of the rational number, for any rational $q \in \mathbb{Q} \cap (0,1]$, $\mu(0,q] = cq$. By continuity, it's easy to see that if we have $A_n = (0,q_n]$ where $q_n$ is a sequence of rationals converging from downwards to a real number, we see that $A_n \uparrow (0,r]$, and 
						\begin{equation*}
							\mu(0,r) = \lim_{n \to \infty} \mu(A_n) = \lim_{n \to \infty} \mu(0,q_n] = c \cdot \lim_{n \to \infty} q_n = cr
						\end{equation*}
						as $q_n \to r$. 
						By translation invariance, we have any interval $(a,b]$ in the reals where $a<b$ has measure, $\mu(a,b]= \mu(0,b-a] = c(b-a)$. It's straightforward to see that $\mu = c\lambda$ on the algebra $\Sigma_0$ as defined in question \ref{q:5}, which is trivially a $\pi$ system, thus they agree on $\mathcal{B}(0,1]$
					\end{proof}
			\end{itemize}
			\item find $\limsup$ and $\liminf$ of $A_n = [\frac{1}{n}, 1]$ if $n$ is odd, and $A_n = [-1,-\frac{1}{n}]$ if $n$ is even. 
			\begin{proof}[Solution.]
				Take any $a \in [-1,1]\setminus\{0\}$, then $a$ is in infinitely many $A_n$'s so we can conclude that $[-1,1]\setminus\{0\} \subseteq \limsup_{n \to \infty} A_n$. Meanwhile, $A_n \subseteq [-1,1]\setminus\{0\}$ for all $n$, so $\limsup_n A_n = [-1,1]\setminus\{0\}$. As for $\liminf$, since $A_n \cap A_{n+1} = \emptyset$, $\liminf_n A_n = \emptyset$. 
			\end{proof}
			
			
			\item Letting $\mathbb{I}_A$ be the indicator function for if some element $s$ is in $A$ or not. Show that $\mathbb{I}_{\limsup_n A_n} (s) = \limsup_n (\mathbb{I}_{A_n}(s))$
				\begin{proof}
					\begin{align*}
						\mathbb{I}_{\limsup_n A_n} (s)=1 &\iff s \in \limsup_n A_n(s)\\
						&\iff \exists \{A_{n_k} \colon k \geq 1\} \subseteq \{A_n\}_{n \geq 1} \text{ s.t } s \in A_{n_k} \forall k\\
						&\iff  \mathbb{I}_{A_{n_k}}(s) = 1 \text{ for all }k\\
						&\iff \limsup_{n \to \infty}(\mathbb{I}_{A_n}(s)) =1
					\end{align*}
				\end{proof}
		\item Given a probability space $(\Omega, \mathcal{F}, \mathbb{P})$. Let $\{A_n\}_{n \geq 1} \subseteq \mathcal{F}$ be a sequence of events. 
			\begin{itemize}
				\item[(i)] Prove that $\lim_{n \to \infty}\mathbb{P}\big(\liminf_{k \to \infty}(A_n\setminus A_k)\big) = 0$
				\begin{proof}[Solution.]
					First of all, for any $n \geq 1$, $\liminf_{k \to \infty} (A_n \setminus A_k)=\emptyset$ since eventually $k=n$. Thus the $\limsup_{n \to \infty}$ of the latter is obviously $\emptyset$. Otherwise, another proof would be to assume that there is an element $w \in \limsup_{n \to \infty} \liminf_{k \to \infty} (A_n \setminus A_k)$, then that means $w$ is in infinitely many $(\liminf_{k \to \infty}(A_n \setminus A_k))_n$. Suppose for the sake of contradiction that $w$ is in one of the latter for some index $n_l$, that means $w \in \liminf_{k \to \infty} (A_{n_l} \setminus A_k)$, that means there exists a $k$ such that for all $j > k$, $w \in A_{n_l}\setminus A_j$, but if $k=n_l$ that's impossible, which results in a contradiction. 
				\end{proof}
				\item[(ii)] Note that $A_n \setminus A^* = A_n \cap (\limsup_{k \to \infty} A_k)^c = A_n \cap \liminf_{k \to \infty} A_k^c = \liminf_{k \to \infty} (A_n \cap A_k^c)$, which by $(i)$, gives us the desired result. Next applying (i) to a sequence $\{A_n^c}\}_{n \geq 1}$, we have that 
				\begin{align*}
					0 = \lim_n \mathbb{P}(\liminf_k (A_n^c \setminus A_k^c))= \lim_n \mathbb{P}(A_* \setminus A_n)
				\end{align*}
				\item[(iii)] Recall that set theoretical symmetric difference is $A \Delta B = (B\setminus A) \cup (A \setminus B)$ 
				. Prove that, assuming $A^* = A_* = A$, then show that $\lim_{n \to \infty}\mathbb{P}(A_n \Delta A)	 =0$.
				\begin{proof}[Solution.]
					First of all, there exists $\mathbb{P}(A_n \Delta A) = \mathbb{P}(A_n \setminus A) + \mathbb{P}(A \setminus A_n)  = \mathbb{P}(A_n \setminus A^*) + \mathbb{P}(A_*\setminus A_n) = 0$ by (ii)
				\end{proof}
				\item[(iv)] Now prove the same thing under a weaker condition 
			\end{itemize}

\end{enumerate}


\end{document}
